\scanpage{16}
Donc on est sous les conditions du N\textsuperscript{o} 1, avec $A_0=A/\unsure{m|\mathfrak m}=k$ un corps. \deleted{Comme} Nous savons par N\textsuperscript{o} 1 prop. 3, cor. 2 qu'on aura
% slanted
\margin{dégager un lemme de réduction}
\diagram{\begin{tikzcd}[ampersand replacement=\&]
A \arrow[r] \arrow[d] \& A' \arrow[r,shift left=.5ex] \arrow[r,shift right=.5ex] \arrow[d] \& A'' \arrow[d] \\
A_0 \arrow[r] \& A'_0 \arrow[r,shift left=.5ex] \arrow[r,shift right=.5ex] \& A''_0
\end{tikzcd}}
avec la deuxième ligne exacte ; notons qu'on aura $A'_0=A'/mA'=A'/m'$, $A'_0$ est \uline{réduit}. De plus,
\deleted{\unsure{quand|quant|posant}}
\[
\begin{aligned}
\Delta_0=\operatorname{Ker}(A''_0\mathrel{\unsure{-|\longrightarrow}}A'_0)&\simeq\operatorname{Ker}(A''/mA''\longrightarrow A'/mA'),\\
&\simeq\operatorname{Ker}(A''/mA''\mathrel{-}A''/\Delta+mA'')\\
&\simeq(\Delta+mA'')/mA''\simeq\Delta/\Delta\cap mA'',
\end{aligned}
\]
et notant que $\Delta\simeq\operatorname{Ker}(A''\mathrel{-}A')\longrightarrow\Delta_0\simeq\operatorname{Ker}(A''_0\mathrel{-}A'_0)$ est \uline{surjectif}, on voit que \deleted{\unsure{$A$|$\Delta$}} l'hypothèse faite pour $A'/A$ implique que pour tt $x'_0\in A'_0-A_0$, on a $\delta_0(x'_0)\in A'_0$ est un \uline{générateur} de $\deleted{\unsure{A''}}\Delta_0$. Utilisant le fait que les composants locaux de $A'_0$ \deleted{\unsure{$=$|\text{et}}} [qui sont des corps] sont des extensions \uline{triviales} \inserted{de $k=A_0$}, \deleted{et} donc $A'_0\simeq\underbrace{k\times\ldots\times k}_{n}$, avec $n\geq2$, et prenant $x'_0=(1,0,\ldots,0)$ on trouve $n=2$, O.K.
% archive sheet number: 15
